%====================================================================
\section{Annexe : choix technique détaillé}
\label{annexe:choix}
%====================================================================

Le tableau~\ref{tab:solutions} compare les solutions envisagées. La machine disponible est un PC Linux (Debian 12, 12 cœurs, 14~Go de RAM, \textbf{sans GPU NVIDIA}) sur lequel Docker est installé.

\begin{table}[H]
\centering
\footnotesize
\begin{tabularx}{\textwidth}{>{\raggedright\arraybackslash}p{2.2cm}>{\raggedright\arraybackslash}X>{\raggedright\arraybackslash}X>{\raggedright\arraybackslash}X>{\raggedright\arraybackslash}X}
\toprule
 & \textbf{AlphaFold Server (AF3)}~\cite{alphafold_server} & \textbf{AlphaFold3 local}~\cite{alphafold3_github} & \textbf{ColabFold local, Docker (AF2)}~\cite{colabfold_github} & \textbf{AlphaFold2 officiel}~\cite{alphafold_github} \\
\midrule
À installer & Rien (navigateur, compte Google) & Docker, pilotes NVIDIA, code DeepMind, poids sur demande & Docker + image ColabFold (\textasciitilde 10~Go) + poids (\textasciitilde 5~Go) & Docker, pilotes NVIDIA, scripts DeepMind \\
Calcul & Serveurs Google & Local & Local (CPU ou GPU) & Local \\
Ressources & Aucune localement & GPU obligatoire (A100/H100 80~Go conseillé ; 16~Go limite à 1\,280 jetons) & CPU suffisant pour $<$ 300 résidus & GPU recommandé, $\geq$ 64~Go RAM \\
Bases de données & Aucune & \textasciitilde 630~Go & Aucune : MSA via serveur MMseqs2 & \textasciitilde 2,6~To, ou \textasciitilde 500~Go en \texttt{reduced\_dbs} \\
Avantages & Immédiat, AF3 complet, gratuit & Tous les ligands, scriptable, pas de quota & Reproductible, scriptable & Pipeline de référence, hors ligne \\
Inconvénients & Quota journalier (30 tâches), pas d'API, liste de ligands restreinte, non commercial & Impossible sans GPU ; licence stricte, poids non redistribuables & AF2 seulement ; lent sur CPU & AF2 seulement ; irréaliste sur la machine disponible \\
\bottomrule
\end{tabularx}
\caption{Comparaison des solutions envisagées.}
\label{tab:solutions}
\end{table}

\paragraph{Solution retenue pour AlphaFold3 : AlphaFold Server.} L'installation locale d'AlphaFold3 exige un GPU récent, absent de la machine disponible : elle est exclue. AlphaFold Server~\cite{alphafold_server} donne accès gratuitement au modèle complet depuis un navigateur, pour un usage non commercial. Son quota (30 tâches par jour, constaté à la connexion) suffit largement pour ce projet. Ses limites sont l'absence d'API, qui empêche l'automatisation, et une liste de ligands plus restreinte que la version locale.

\paragraph{Référence AlphaFold2 : ColabFold~1.5.5 dans Docker.} Pour comparer, le déploiement d'AlphaFold2 a été conservé. ColabFold~\cite{mirdita2022} utilise \emph{les mêmes réseaux et les mêmes poids} qu'AlphaFold2 ; il remplace seulement la construction du MSA (2,6~To de bases) par une recherche MMseqs2~\cite{steinegger2017} sur un serveur public. Il tourne sur CPU, dans un environnement figé par Docker.

\paragraph{AlphaFold2 officiel en mode \texttt{reduced\_dbs}.} Il est aussi possible de déployer l'installation officielle de DeepMind~\cite{alphafold_github}. Le disque ne permet pas d'accueillir les \textasciitilde 2,6~To de la configuration complète : il faut donc utiliser l'option \texttt{--db\_preset=reduced\_dbs}. Elle ne modifie que la construction du MSA ; le réseau, ses poids et la prédiction restent identiques. En mode complet, AlphaFold interroge UniRef90 et MGnify avec JackHMMER, puis BFD (\textasciitilde 1,8~To, en partie métagénomique) et UniRef30 (\textasciitilde 200~Go) avec HHblits. En mode \texttt{reduced\_dbs}, ces deux grosses bases sont remplacées par \textbf{small BFD} (\textasciitilde 17~Go), un sous-ensemble de BFD débarrassé de la plupart de ses séquences redondantes. L'alignement est un peu moins riche, ce qui pénalise surtout les protéines \og orphelines \fg{} ; pour la grande majorité des protéines, la différence est faible, et la recherche est plus rapide.

%====================================================================
\section{Annexe : mise en œuvre détaillée}
\label{annexe:miseenoeuvre}
%====================================================================

\subsection{Prérequis}

Le tableau~\ref{tab:prerequis} résume ce dont il faut disposer avant de commencer. Les versions indiquées sont celles qui ont été testées.

\begin{table}[H]
\centering
\footnotesize
\begin{tabularx}{\textwidth}{>{\raggedright\arraybackslash}p{2.4cm}>{\raggedright\arraybackslash}X>{\raggedright\arraybackslash}X}
\toprule
 & \textbf{AlphaFold3 (AlphaFold Server)} & \textbf{AlphaFold2 (ColabFold dans Docker)} \\
\midrule
Système & N'importe lequel, avec un navigateur récent & Linux (testé : Debian~12, noyau 6.12), processeur x86\_64 \\
Compte & Compte Google & Aucun \\
Logiciels & Aucun & Docker (testé : 29.8.1), utilisable sans \texttt{sudo} (groupe \texttt{docker}) ; Python~3.11 avec numpy et matplotlib pour l'analyse \\
Matériel & Aucun & CPU suffisant (testé : 12 cœurs, 14~Go de RAM) ; GPU NVIDIA facultatif \\
Disque & Aucun & \textasciitilde 20~Go (image Docker \textasciitilde 10~Go, poids 5,3~Go) \\
Réseau & Connexion Internet & Connexion Internet : téléchargements initiaux, puis MSA calculé par le serveur \texttt{api.colabfold.com} à chaque prédiction \\
\bottomrule
\end{tabularx}
\caption{Prérequis des deux solutions déployées.}
\label{tab:prerequis}
\end{table}

\subsection{AlphaFold3 via AlphaFold Server}
\label{sec:af3_server}

\paragraph{Installation.} Aucune : il suffit de se connecter sur \url{https://alphafoldserver.com/} avec un compte Google et d'accepter les conditions d'utilisation. L'interface est en anglais ; le compteur \og Remaining jobs \fg{} indique le quota du jour (30).

\paragraph{Soumettre une nouvelle séquence.}
\begin{enumerate}
  \item Laisser \og Entity type \fg{} sur \emph{Protein} et \og Copies \fg{} sur 1, puis coller la séquence dans le champ \og Paste sequence or fasta \fg{} (séquence brute ou format FASTA).
  \item Pour un complexe, cliquer sur \og Add entity \fg{} et choisir une autre chaîne, de l'ADN, de l'ARN, un ligand ou un ion : c'est là que se voit la différence avec AlphaFold2.
  \item Cliquer sur \og Continue and preview job \fg{}, nommer la tâche, puis cliquer sur \og Confirm and submit job \fg{}.
  \item La tâche apparaît dans l'historique avec une icône de chargement (filtre \og In progress \fg{}), puis passe dans \og Completed \fg{} quand elle est terminée (quelques minutes pour une petite protéine).
\end{enumerate}

\paragraph{Résultats.} Le résultat s'affiche directement dans le navigateur : structure 3D colorée par pLDDT, matrice PAE, pTM et ipTM. Le bouton \og Download \fg{}, en haut de la page, fournit une archive ZIP. Ses fichiers principaux sont les suivants (\texttt{<nom>} est le nom donné à la tâche) :
\begin{itemize}
  \item \texttt{fold\_<nom>\_model\_0.cif} à \texttt{model\_4.cif} : les 5 structures prédites (format mmCIF), le modèle~0 étant le mieux classé ;
  \item \texttt{fold\_<nom>\_summary\_confidences\_*.json} : les scores globaux de chaque modèle (pTM, ipTM, score de classement, fraction désordonnée, présence de collisions) ;
  \item \texttt{fold\_<nom>\_full\_data\_*.json} : les scores détaillés (pLDDT par atome, matrice PAE, probabilités de contact) ;
  \item \texttt{fold\_<nom>\_job\_request.json} : la requête soumise, utile pour la traçabilité et pour relancer la même tâche ;
  \item \texttt{msas/} et \texttt{templates/} : les alignements et les gabarits de la PDB utilisés par le modèle.
\end{itemize}

\paragraph{Difficultés et points d'attention.}
\begin{itemize}
  \item \textbf{Pas d'automatisation} : aucune API n'est proposée ; chaque tâche se soumet à la main, ce qui rend les séries de prédictions fastidieuses.
  \item \textbf{Quota} : 30 tâches par jour. Il faut donc préparer les soumissions plutôt que procéder par essais successifs.
  \item \textbf{Ligands limités} : seule une liste restreinte de ligands et d'ions est proposée ; un médicament quelconque (formule SMILES) n'est accepté que par la version locale.
\end{itemize}

\subsection{AlphaFold2 via ColabFold (référence)}

\paragraph{Installation, dans l'ordre.} Les commandes exactes figurent dans le \texttt{README.md} du dépôt :
\begin{enumerate}
  \item vérifier que Docker fonctionne sans \texttt{sudo} : \texttt{docker run --rm hello-world} ;
  \item télécharger l'image, dans une version figée :\\
  \texttt{docker pull ghcr.io/sokrypton/colabfold:1.5.5-cuda12.2.2} ;
  \item télécharger les poids d'AlphaFold2 (5,3~Go, une seule fois) avec \texttt{python -m colabfold.download}, exécuté dans le conteneur, vers un dossier \texttt{cache/} monté en volume ; \texttt{scripts/predict.sh} le fait automatiquement s'ils sont absents ;
  \item installer les dépendances Python de l'analyse : \texttt{pip install -r scripts/requirements.txt}.
\end{enumerate}

\paragraph{Soumettre une nouvelle séquence.}
\begin{enumerate}
  \item Créer un fichier FASTA dans \texttt{data/}, avec \textbf{un en-tête court} et une seule séquence protéique. Par exemple, depuis UniProt :\\
  \texttt{curl -s https://rest.uniprot.org/uniprotkb/P0DP23.fasta \textbackslash}\\
  \texttt{\ \ | sed '1s/.*/>P0DP23/' > data/P0DP23.fasta}
  \item Lancer la prédiction : \texttt{scripts/predict.sh data/P0DP23.fasta}. Le script exécute \texttt{colabfold\_batch --num-models 5} dans le conteneur : MSA via MMseqs2, puis prédiction par les 5 modèles AlphaFold2. Il ajoute \texttt{--gpus all} si un GPU est détecté.
  \item Analyser : \texttt{python scripts/analyse.py results/P0DP23 -o figures/P0DP23}, qui extrait le pLDDT et la PAE et produit les figures dans \texttt{figures/P0DP23/}.
\end{enumerate}

\paragraph{Résultats.} \sloppy Les résultats sont écrits dans \texttt{results/<nom>/} et le journal complet, avec la durée, dans \texttt{logs/<nom>.log}. Les fichiers importants sont décrits au tableau~\ref{tab:fichiers} (annexe~\ref{annexe:calmoduline}).\par\fussy

\paragraph{Difficultés rencontrées et solutions.}
\begin{itemize}
  \item \textbf{Pas de GPU} : JAX bascule automatiquement sur le CPU. Le calcul fonctionne, mais il est lent : 1~h~57 pour 149 résidus et 5 modèles.
  \item \textbf{Fichiers appartenant à \texttt{root}} : Docker écrit en \texttt{root} par défaut ; il faut le lancer avec \texttt{--user \$(id -u):\$(id -g) -e HOME=/tmp} (déjà fait dans \texttt{predict.sh}).
  \item \textbf{Noms de fichiers très longs} : ColabFold nomme les sorties d'après l'en-tête FASTA complet ; il faut le raccourcir (\texttt{>P0DP23}).
  \item \textbf{Mémoire} : les avertissements \texttt{Allocation exceeds 10\% of free system memory} sont sans conséquence ; la consommation plafonne à environ 4~Go.
  \item \textbf{Accès à Docker refusé} (\texttt{permission denied ... docker.sock}) : l'utilisateur doit être ajouté au groupe \texttt{docker} (\texttt{sudo usermod -aG docker \$USER}), puis se reconnecter.
\end{itemize}

\subsection{AlphaFold2 officiel en mode \texttt{reduced\_dbs}}

Cette installation n'est réaliste que sur une machine plus puissante (GPU recommandé, $\geq$ 64~Go de RAM, \textasciitilde 500~Go de disque). Les étapes, dans l'ordre, sont celles du dépôt de DeepMind~\cite{alphafold_github} :
\begin{enumerate}
  \item cloner le dépôt : \texttt{git clone https://github.com/google-deepmind/alphafold.git} ;
  \item télécharger les bases réduites et les poids (nécessite \texttt{aria2c}) :\\
  \texttt{scripts/download\_all\_data.sh <DOSSIER\_BASES> reduced\_dbs} ;
  \item construire l'image Docker : \texttt{docker build -f docker/Dockerfile -t alphafold .} ;
  \item installer le lanceur : \texttt{pip install -r docker/requirements.txt} ;
  \item lancer une prédiction :\\
  \texttt{python docker/run\_docker.py --fasta\_paths=séquence.fasta --db\_preset=reduced\_dbs \textbackslash}\\
  \texttt{\ \ --data\_dir=<DOSSIER\_BASES> --output\_dir=<SORTIE> --max\_template\_date=2022-01-01}\\
  (ajouter \texttt{--use\_gpu=false} sur une machine sans GPU).
\end{enumerate}
Les résultats sont écrits dans \texttt{<SORTIE>/<nom>/} : \texttt{ranked\_0.pdb} est la meilleure structure, \texttt{ranking\_debug.json} donne le classement des 5 modèles, les fichiers \texttt{result\_model\_*.pkl} contiennent les scores (pLDDT, PAE), et le dossier \texttt{msas/} les alignements. La principale difficulté est l'espace disque : la configuration complète (\textasciitilde 2,6~To) ne tient pas sur la machine disponible, d'où le mode \texttt{reduced\_dbs}.

\subsection{Limites de la solution déployée}

La solution fonctionne dans les conditions du projet, mais plusieurs situations la mettraient en défaut :
\begin{itemize}
  \item \textbf{Serveurs externes indisponibles.} ColabFold dépend du serveur public \texttt{api.colabfold.com} pour le MSA, et AlphaFold3 entièrement d'AlphaFold Server. Si l'un d'eux est arrêté, saturé ou limité, la prédiction devient impossible. Pour ColabFold, il faudrait alors héberger ses propres bases (\texttt{colabfold\_search}, \textasciitilde 1~To). Pour AlphaFold Server, le risque est réel : l'équipe AlphaFold de DeepMind a été dissoute en juillet 2026~\cite{gigazine2026}, et rien ne garantit que le service restera gratuit ou ouvert.
  \item \textbf{Protéines longues ou complexes volumineux.} Sur CPU, le temps de calcul devient prohibitif au-delà d'environ 400 résidus ; il faut alors un GPU ou Google Colab. AlphaFold Server limite aussi la taille totale d'une tâche (5\,000 jetons).
  \item \textbf{Séquences confidentielles.} Les deux solutions envoient la séquence à un serveur externe (MSA pour ColabFold, tout le calcul pour AlphaFold Server). Elles ne conviennent donc pas à des séquences non publiées ou protégées.
  \item \textbf{Usage commercial.} AlphaFold Server est réservé à un usage non commercial ; une entreprise ne peut pas l'utiliser pour ses projets. AlphaFold2 (ColabFold), sous licence libre, n'a pas cette restriction.
  \item \textbf{Grand nombre de prédictions.} Sans API et avec 30 tâches par jour, AlphaFold Server ne permet pas de criblage à grande échelle.
  \item \textbf{Évolution de l'environnement.} L'image Docker de ColabFold est construite pour x86\_64 et n'a pas été testée sur processeur ARM (Mac M1/M2). Si l'image ou les poids venaient à être retirés de leurs dépôts, la version figée ne pourrait plus être réinstallée.
\end{itemize}
